% part: intuitionistic-logic
% chapter: propositions-as-types
% section: proof-terms

\documentclass[../../../include/open-logic-section]{subfiles}

\begin{document}

\olfileid[hi]{int}{pty}{ter}

\olsection{प्रमाण पद}

\Log{N1i} में हम मान्यताओं को चरों (जैसे $!A^x$) से मुक्तिकरण चिह्न देते हैं।
\Log{N2i} में प्रत्येक अनुक्रम $\Gamma \Sequent !A$ के बाएँ प्रसंग~$\Gamma$
में चर--!!{formula} युग्मों की सूची होती है। अतः \Log{N2i} में किसी !!{proof}
का प्रत्येक अनुक्रम न केवल यह सूचना रखता है कि अब तक !!{proof} ने क्या
!!{prove} किया है (अर्थात्~$!A$), बल्कि यह भी कि प्रत्येक बिंदु पर कौन-सी
मान्यताएँ खुली हैं (अर्थात्~$\Gamma$)। यदि हर अनुक्रम में यह सूचना भी शामिल
करें कि $!A$ को \emph{कैसे} !!{prove} किया गया? इसके लिए हम \emph{पद-चिह्नित}
प्रणाली \Log{tN2i} परिभाषित कर सकते हैं, जिसमें अनुक्रमों का रूप $\Gamma
\Sequent N : !A$ है। $\Gamma$ और $!A$ पहले जैसे हैं: $x: !B$ रूप के युग्मों
का समुच्चय और वह !!{formula} जिसे अनुक्रम पर समाप्त होने वाले उप-!!{proof} ने
~$\Gamma$ की मान्यताओं से निगमित दिखाया है। $N$ ऐसा पद होगा जो अनुक्रम पर
समाप्त होने वाले !!{proof} को संहिताबद्ध करता है।

हम प्रत्येक अनुक्रम को जो पद देते हैं वे चर $x$, $y$, \dots से ऐसे फलन-चिह्नों
द्वारा बनते हैं जो अब तक किए गए अनुमानों को संहिताबद्ध करते हैं। प्रत्येक
अनुमान-नियम के लिए अलग फलन-चिह्न चाहिए। !!{introduction} नियमों से संबंधित
फलन-चिह्न ``रचयिता'' और !!{elimination} नियमों से संबंधित फलन-चिह्न
``विघटक'' कहलाते हैं। प्रत्येक फलन-चिह्न उतने तर्क लेता है जितने नियम के
आधार-वाक्य हैं। उदाहरणतः \Intro{\land} और \Elim{\land}[1] के फलन-चिह्न
$c^\land$ (द्विस्थानी) और $d^\land_1$ (एकस्थानी) हो सकते हैं। तब \Log{tN2i}
में ये नियम बनते हैं
\[
\Axiom$\Gamma_1 \fCenter N: !A$
\Axiom$\Gamma_2 \fCenter M: !B$
\RightLabel{\Intro{\land}}
\BinaryInf$\Gamma_1, \Gamma_2 \fCenter c^\land(N, m) : !A \land !B $
\DisplayProof
\quad
\Axiom$\Gamma \fCenter N: !A \land !B$
\RightLabel{\Elim{\land}[1]}
\UnaryInf$\Gamma \fCenter d^\land_1(N) : !A$
\DisplayProof
\]
यदि कोई अनुमान-नियम किसी मान्यता को मुक्त करता है—अर्थात् नियम के किसी
आधार-वाक्य में ऐसा चिह्नित !!{formula}~$x:!A$ है जो निष्कर्ष के प्रसंग में नहीं
है—तो प्रमाण पद में भी इसे बताना होगा। ऐसे नियमों से संबंधित फलन-चिह्न संगत
चरों को \emph{आबद्ध} करते हैं। उदाहरणतः \Intro{\lif} नियम आधार-वाक्य $x:!A,
\Gamma \Sequent !B$ से $\Gamma \Sequent !A \lif !B$ तक जाता है। रचयिता
$c^{\lif}$ एक तर्क लेता है, पर उसे यह भी बताना होगा कि चर $x$ आबद्ध हो जाता
है। यह बताने का प्रमाण पद का ढंग है कि~$x$ चिह्न वाली मान्यता मुक्त हो गई है।
उदाहरणतः यदि आधार-वाक्य का प्रमाण पद~$N$ है, तो निष्कर्ष का प्रमाण पद
$c^{\lif}([x]N)$ और नियम इस प्रकार लिख सकते हैं
\[
\Axiom$x:!A, \Gamma \fCenter N: !B$
\RightLabel{\Intro{\lif}}
\UnaryInf$\Gamma \fCenter c^{\lif}([x]N) : !A \lif !B$
\DisplayProof
\]

किसी !!{proof} का पूर्ण पुनर्निर्माण करने के लिए कुछ अतिरिक्त सूचना चाहिए।
यदि किसी नियम के आधार-वाक्य के !!{formula}s निष्कर्ष के !!{formula} को पूरी
तरह निर्धारित नहीं करते, तो वह सूचना पद में जोड़नी होगी। \Intro{\lif} में ऐसा
ही है, इसलिए हम $c^{\lif}([x:!A]N)$ लिखेंगे। \Intro{\lor} और~\FalseInt में भी
ऐसा है, इसलिए उस सूचना को रखने वाले फलन-चिह्न प्रयुक्त करते हैं, उदाहरणतः
\[
\Axiom$\Gamma \fCenter N:!A$
\RightLabel{\Intro{\lor}[1]}
\UnaryInf$\Gamma \fCenter c^{\lor{!B}}_1:!A \lor !B$
\DisplayProof
\quad
\Axiom$\Gamma \fCenter N:\lfalse$
\RightLabel{\FalseInt}
\UnaryInf$\Gamma \fCenter d^\lfalse_{!A}(N):!A$
\DisplayProof
\]

इन विचारों से ऐसी अनुक्रम-शैली प्राकृतिक निगमन प्रणाली बन सकती है जिसमें हर
अनुक्रम $\Gamma \Sequent N : !A$ रूप का है, जहाँ $N$ एक \emph{प्रमाण पद} है।

ऐसे प्रमाण पदों और तथाकथित \emph{प्रकारित लैम्ब्डा कलन} के पदों में निकट संबंध
है। इस संबंध के अनुरूप हम ऊपर के सामान्य रचयिता और विघटक चिह्नों के बजाय
लैम्ब्डा कलन के साहित्य में प्रयुक्त चिह्न लेंगे। उदाहरणतः
$c^{\lif}([x:!A]N)$ के बजाय $\lambd[\typeof{x}{!A}][N]$, $d^{\lif}(N,M)$ के
बजाय $(NM)$, $c^\land(N,M)$ के बजाय $\pair{N}{M}$, और~$d^\land_i(N)$ के बजाय
$\proj{i}{N}$ लिखते हैं।

\begin{defn}
  प्रमाण पद निम्न नियमों द्वारा आगमनिक रूप से उत्पन्न होते हैं:
  \begin{enumerate}
  \item प्रत्येक चर $x$ प्रमाण पद है (अभिगृहीत)।
  \item यदि $x$ चर, $!A$ कोई !!a{formula}, और $N$ प्रमाण पद है, तो
    $\lambd[\typeof{x}{!A}][N]$ भी प्रमाण पद है~(\Intro\lif)।
  \item यदि $N$ और $M$ प्रमाण पद हैं, तो $(NM)$ भी प्रमाण पद है~(\Elim\lif)।
  \item यदि $N$ और $M$ प्रमाण पद हैं, तो $\pair{N}{M}$ प्रमाण पद है
    ~(\Intro\land)।
  \item यदि $N$ प्रमाण पद है, तो $\proj{i}{N}$ भी प्रमाण पद है, जहाँ $i$, $1$
    या~$2$ है~(\Elim\land[i])।
  \item यदि $N$ प्रमाण पद और $!A$ कोई सूत्र है, तो $\inj[!A]{i}{!M}$ प्रमाण
    पद है, जहाँ $i$, $1$ या~$2$ है~(\Intro\lor[i])।
  \item यदि $M$, $N_1$, $N_2$ प्रमाण पद और $x_1$, $x_2$ चर हैं, तो
    $\dcase{M}{x_1}{N_1}{x_2}{N_2}$ प्रमाण पद है~(\Elim\lor)।
  \item यदि $N$ प्रमाण पद और $!A$ कोई !!a{formula} है, तो $\abort{!A}{N}$
    प्रमाण पद है~(\FalseInt)।
  \end{enumerate}
\end{defn}

हर प्रमाण पद किसी !!a{proof} का प्रमाण पद नहीं है। उदाहरणतः पद
$\pair{N}{M}$,~\Intro{\land} अनुमान पर समाप्त होने वाले प्रमाण को संहिताबद्ध
करता है, इसलिए उसका अंतिम !!{formula} कोई संयोजन~$!A \land !B$ है। यह
~\Elim{\lif} नियम का मुख्य आधार-वाक्य नहीं हो सकता, इसलिए $(\pair{N}{M}P)$
सही प्रमाण का प्रमाण पद नहीं हो सकता। केवल \Log{tN2ip} के नियमों का अनुसरण
करने वाले प्रमाण पद \emph{सही} प्रमाण पद हैं। ये नियम \olref{tab:tN2ip} में
दिए गए हैं।

\subfile{rules-tN2}

\end{document}
